\documentclass[11pt,a4paper]{article}
\usepackage{fontspec,amsmath,amssymb,amsthm,microtype}
\setmainfont{lmroman10-regular.otf}[BoldFont=lmroman10-bold.otf,ItalicFont=lmroman10-italic.otf,BoldItalicFont=lmroman10-bolditalic.otf]
\usepackage{xcolor}
\usepackage[most]{tcolorbox}
\usepackage{mathrsfs}
\providecommand{\lt}{<}
\providecommand{\gt}{>}
\setmonofont{DejaVuSansMono.ttf}[BoldFont=DejaVuSansMono-Bold.ttf,ItalicFont=DejaVuSansMono-Oblique.ttf,BoldItalicFont=DejaVuSansMono-BoldOblique.ttf,Scale=0.86]
\newfontfamily\leanfallback{STIXTwoMath-Regular.otf}[Scale=0.86]
\definecolor{leanrule}{gray}{0.62}
\definecolor{leanback}{gray}{0.962}
\definecolor{leantitle}{gray}{0.22}
\definecolor{leanmuted}{gray}{0.38}
\newtcolorbox{leanbox}[1]{enhanced,breakable,colback=leanback,colframe=leanrule,
  boxrule=0.4pt,arc=0.6mm,left=2.4mm,right=2.4mm,top=1.8mm,bottom=1.8mm,
  before skip=0.8em,after skip=1em,fonttitle=\small\bfseries,coltitle=leantitle,
  colbacktitle=leanback,titlerule=0.3pt,title={#1}}
\newlength{\leancharwidth}
\newenvironment{leancode}{\par\smallskip\begingroup\ttfamily\footnotesize\raggedright
  \settowidth{\leancharwidth}{M}\setlength{\parskip}{0pt}\setlength{\parindent}{0pt}
  \frenchspacing}{\par\endgroup\smallskip}
\newcommand{\leanline}[2]{\hangindent=\dimexpr#1\leancharwidth+4\leancharwidth\relax
  \hangafter=1\noindent\hspace*{#1\leancharwidth}#2\par}
\newcommand{\leansource}[1]{\par{\footnotesize\color{leanmuted}#1}\par}
\newcommand{\leanlabel}[1]{\par\smallskip\noindent{\small\bfseries #1}\par}
\usepackage[margin=25mm]{geometry}
\usepackage[unicode,hidelinks]{hyperref}
\hypersetup{pdfauthor={LangevinLean Project}}
\setlength{\emergencystretch}{3em}
\setlength{\parskip}{0.4em}
\setlength{\parindent}{0pt}
\allowdisplaybreaks[2]
\raggedbottom
\title{\texorpdfstring{Actual stationary-pair transition operators}{Actual stationary-pair transition operators}}
\author{LangevinLean Project}
\date{}
\begin{document}\maketitle
Draft mathematical exposition; these notes may contain mistakes.


For a stationary pair law \(J\) with both marginals \(\mu\), pullback by the first and second coordinates gives actual isometries \(U_0,U_1:L^2(\mu;F)\to L^2(J;F)\). The adjoint compression \(P=U_0^*U_1\) is a genuine contraction. Reversal symmetry of the pair law makes it self-adjoint.


\section{\texorpdfstring{Constructing the contraction from the two actual marginals}{Constructing the contraction from the two actual marginals}}


\textbf{Claim.} The coordinate evaluations preserve norms and

\[\langle Pf,g\rangle=\langle U_1f,U_0g\rangle,\qquad\|Pf\|\le\|f\|.\]

\textbf{Proof.} Each marginal equality makes the corresponding coordinate map measure-preserving, hence its pullback is an isometry. Define the transition by adjoint compression. The adjoint identity proves the pairing formula. The first evaluation and its adjoint have norm at most one, while the second evaluation preserves norms, giving the pointwise and operator contraction bounds.

\begin{leanbox}{Lean: \texttt{evalFst}, \texttt{evalSnd}, \texttt{transition} and 3 more}
\leansource{\texttt{evalFst} \quad WassersteinMixing/\allowbreak{}Semigroup.lean, lines 28–29}
\begin{leancode}
\leanline{0}{def\ evalFst\ :\ Lp\ F\ 2\ μ\ →ₗᵢ[{\leanfallback 𝕜}]\ Lp\ F\ 2\ J\ :=}
\leanline{2}{Lp.compMeasurePreservingₗᵢ\ {\leanfallback 𝕜}\ Prod.fst\ ⟨measurable\_fst,\ hfst⟩}
\end{leancode}
\leansource{\texttt{evalSnd} \quad WassersteinMixing/\allowbreak{}Semigroup.lean, lines 32–33}
\begin{leancode}
\leanline{0}{def\ evalSnd\ :\ Lp\ F\ 2\ μ\ →ₗᵢ[{\leanfallback 𝕜}]\ Lp\ F\ 2\ J\ :=}
\leanline{2}{Lp.compMeasurePreservingₗᵢ\ {\leanfallback 𝕜}\ Prod.snd\ ⟨measurable\_snd,\ hsnd⟩}
\end{leancode}
\leansource{\texttt{transition} \quad WassersteinMixing/\allowbreak{}Semigroup.lean, lines 37–39}
\begin{leancode}
\leanline{0}{def\ transition\ :\ Lp\ F\ 2\ μ\ →L[{\leanfallback 𝕜}]\ Lp\ F\ 2\ μ\ :=}
\leanline{2}{(evalFst\ ({\leanfallback 𝕜}\ :=\ {\leanfallback 𝕜})\ μ\ J\ hfst).toContinuousLinearMap.adjoint\ ∘L}
\leanline{4}{(evalSnd\ ({\leanfallback 𝕜}\ :=\ {\leanfallback 𝕜})\ μ\ J\ hsnd).toContinuousLinearMap}
\end{leancode}
\leansource{\texttt{transition\_\allowbreak{}inner} \quad WassersteinMixing/\allowbreak{}Semigroup.lean, lines 42–46}
\begin{leancode}
\leanline{0}{theorem\ transition\_inner\ (f\ g\ :\ Lp\ F\ 2\ μ)\ :}
\leanline{4}{inner\ {\leanfallback 𝕜}\ (transition\ ({\leanfallback 𝕜}\ :=\ {\leanfallback 𝕜})\ μ\ J\ hfst\ hsnd\ f)\ g\ =}
\leanline{6}{inner\ {\leanfallback 𝕜}\ (evalSnd\ ({\leanfallback 𝕜}\ :=\ {\leanfallback 𝕜})\ μ\ J\ hsnd\ f)\ (evalFst\ ({\leanfallback 𝕜}\ :=\ {\leanfallback 𝕜})\ μ\ J\ hfst\ g)}
\end{leancode}
\leansource{\texttt{transition\_\allowbreak{}norm\_\allowbreak{}le} \quad WassersteinMixing/\allowbreak{}Semigroup.lean, lines 49–59}
\begin{leancode}
\leanline{0}{theorem\ transition\_norm\_le\ (f\ :\ Lp\ F\ 2\ μ)\ :}
\leanline{4}{‖transition\ ({\leanfallback 𝕜}\ :=\ {\leanfallback 𝕜})\ μ\ J\ hfst\ hsnd\ f‖\ ≤\ ‖f‖}
\end{leancode}
\leansource{\texttt{transition\_\allowbreak{}opNorm\_\allowbreak{}le} \quad WassersteinMixing/\allowbreak{}Semigroup.lean, lines 62–65}
\begin{leancode}
\leanline{0}{theorem\ transition\_opNorm\_le\ :\ ‖transition\ ({\leanfallback 𝕜}\ :=\ {\leanfallback 𝕜})\ (F\ :=\ F)\ μ\ J\ hfst\ hsnd‖\ ≤\ 1}
\end{leancode}
\leanlabel{How the Lean proof follows the argument}
\texttt{eval\allowbreak{}Fst} and \texttt{eval\allowbreak{}Snd} are actual measure-preserving pullbacks. \texttt{transition} is their continuous adjoint composition. The correlation and norm estimates use the defining adjoint identity and the evaluation isometries.
\end{leanbox}

\section{\texorpdfstring{Reversal symmetry proves self-adjointness}{Reversal symmetry proves self-adjointness}}


\textbf{Claim.} If swapping the two coordinates preserves \(J\), then \(\langle Pf,g\rangle=\langle f,Pg\rangle\), so \(P\) is self-adjoint.



\textbf{Proof.} Pullback by the actual swap map is an isometry on the joint space and exchanges the two evaluation maps. Apply preservation of the inner product by this isometry to the correlation identity, then use the adjoint identity in the other direction. This is the symmetric pairing characterization of a bounded self-adjoint operator.

\begin{leanbox}{Lean: \texttt{reverse}, \texttt{reverse\_\allowbreak{}evalFst}, \texttt{reverse\_\allowbreak{}evalSnd} and 2 more}
\leansource{\texttt{reverse} \quad WassersteinMixing/\allowbreak{}Semigroup.lean, lines 70–71}
\begin{leancode}
\leanline{0}{def\ reverse\ :\ Lp\ F\ 2\ J\ →ₗᵢ[{\leanfallback 𝕜}]\ Lp\ F\ 2\ J\ :=}
\leanline{2}{Lp.compMeasurePreservingₗᵢ\ {\leanfallback 𝕜}\ Prod.swap\ ⟨measurable\_swap,\ hswap⟩}
\end{leancode}
\leansource{\texttt{reverse\_\allowbreak{}evalFst} \quad WassersteinMixing/\allowbreak{}Semigroup.lean, lines 74–86}
\begin{leancode}
\leanline{0}{theorem\ reverse\_evalFst\ (f\ :\ Lp\ F\ 2\ μ)\ :}
\leanline{4}{reverse\ ({\leanfallback 𝕜}\ :=\ {\leanfallback 𝕜})\ J\ hswap\ (evalFst\ ({\leanfallback 𝕜}\ :=\ {\leanfallback 𝕜})\ μ\ J\ hfst\ f)\ =\ evalSnd\ ({\leanfallback 𝕜}\ :=\ {\leanfallback 𝕜})\ μ\ J\ hsnd\ f}
\end{leancode}
\leansource{\texttt{reverse\_\allowbreak{}evalSnd} \quad WassersteinMixing/\allowbreak{}Semigroup.lean, lines 89–103}
\begin{leancode}
\leanline{0}{theorem\ reverse\_evalSnd\ (f\ :\ Lp\ F\ 2\ μ)\ :}
\leanline{4}{reverse\ ({\leanfallback 𝕜}\ :=\ {\leanfallback 𝕜})\ J\ hswap\ (evalSnd\ ({\leanfallback 𝕜}\ :=\ {\leanfallback 𝕜})\ μ\ J\ hsnd\ f)\ =\ evalFst\ ({\leanfallback 𝕜}\ :=\ {\leanfallback 𝕜})\ μ\ J\ hfst\ f}
\end{leancode}
\leansource{\texttt{transition\_\allowbreak{}symmetric} \quad WassersteinMixing/\allowbreak{}Semigroup.lean, lines 105–117}
\begin{leancode}
\leanline{0}{theorem\ transition\_symmetric\ (f\ g\ :\ Lp\ F\ 2\ μ)\ :}
\leanline{4}{inner\ {\leanfallback 𝕜}\ (transition\ ({\leanfallback 𝕜}\ :=\ {\leanfallback 𝕜})\ μ\ J\ hfst\ hsnd\ f)\ g\ =}
\leanline{6}{inner\ {\leanfallback 𝕜}\ f\ (transition\ ({\leanfallback 𝕜}\ :=\ {\leanfallback 𝕜})\ μ\ J\ hfst\ hsnd\ g)}
\end{leancode}
\leansource{\texttt{transition\_\allowbreak{}isSelfAdjoint} \quad WassersteinMixing/\allowbreak{}Semigroup.lean, lines 119–121}
\begin{leancode}
\leanline{0}{theorem\ transition\_isSelfAdjoint\ :\ IsSelfAdjoint\ (transition\ ({\leanfallback 𝕜}\ :=\ {\leanfallback 𝕜})\ (F\ :=\ F)\ μ\ J\ hfst\ hsnd)}
\end{leancode}
\leanlabel{How the Lean proof follows the argument}
\texttt{reverse\_\allowbreak{}eval\allowbreak{}Fst} and \texttt{reverse\_\allowbreak{}eval\allowbreak{}Snd} compare almost-everywhere pullback representatives. \texttt{transition\_\allowbreak{}symmetric} uses the actual reversal isometry and \texttt{transition\_\allowbreak{}is\allowbreak{}Self\allowbreak{}Adjoint} applies the symmetric-operator characterization.
\end{leanbox}

\section{\texorpdfstring{Compression on any common observation space}{Compression on any common observation space}}


\textbf{Claim.} The same construction works for any two isometries into a common Hilbert space: \(P=U_0^*U_1\) has the same pairing formula and norm bound. For equal observations, \(U_0^*U_0=I\).



\textbf{Proof.} The contraction argument uses only the two isometry properties. At equal observations, the adjoint composition of an isometry with itself is exactly the identity.

\begin{leanbox}{Lean: \texttt{compression}, \texttt{compression\_\allowbreak{}inner}, \texttt{compression\_\allowbreak{}norm\_\allowbreak{}le} and 1 more}
\leansource{\texttt{compression} \quad WassersteinMixing/\allowbreak{}Semigroup.lean, lines 132–133}
\begin{leancode}
\leanline{0}{def\ compression\ (initial\ current\ :\ A\ →ₗᵢ[{\leanfallback 𝕜}]\ B)\ :\ A\ →L[{\leanfallback 𝕜}]\ A\ :=}
\leanline{2}{initial.toContinuousLinearMap.adjoint\ ∘L\ current.toContinuousLinearMap}
\end{leancode}
\leansource{\texttt{compression\_\allowbreak{}inner} \quad WassersteinMixing/\allowbreak{}Semigroup.lean, lines 135–138}
\begin{leancode}
\leanline{0}{theorem\ compression\_inner\ (initial\ current\ :\ A\ →ₗᵢ[{\leanfallback 𝕜}]\ B)\ (f\ g\ :\ A)\ :}
\leanline{4}{inner\ {\leanfallback 𝕜}\ (compression\ initial\ current\ f)\ g\ =\ inner\ {\leanfallback 𝕜}\ (current\ f)\ (initial\ g)}
\end{leancode}
\leansource{\texttt{compression\_\allowbreak{}norm\_\allowbreak{}le} \quad WassersteinMixing/\allowbreak{}Semigroup.lean, lines 141–150}
\begin{leancode}
\leanline{0}{theorem\ compression\_norm\_le\ (initial\ current\ :\ A\ →ₗᵢ[{\leanfallback 𝕜}]\ B)\ (f\ :\ A)\ :}
\leanline{4}{‖compression\ initial\ current\ f‖\ ≤\ ‖f‖}
\end{leancode}
\leansource{\texttt{compression\_\allowbreak{}self} \quad WassersteinMixing/\allowbreak{}Semigroup.lean, lines 153–155}
\begin{leancode}
\leanline{0}{theorem\ compression\_self\ (initial\ :\ A\ →ₗᵢ[{\leanfallback 𝕜}]\ B)\ :}
\leanline{4}{compression\ initial\ initial\ =\ ContinuousLinearMap.id\ {\leanfallback 𝕜}\ A}
\end{leancode}
\leanlabel{How the Lean proof follows the argument}
\texttt{compression} abstracts only the isometric evaluation algebra. Its correlation and contraction lemmas are used by the actual Brownian observations; \texttt{compression\_\allowbreak{}self} supplies their exact time-zero identity.
\end{leanbox}

\end{document}
